\section{Limitations}
Our work has several limitations:

\noindent\textbf{Assessment in QA settings.} Our evaluation follows the standard practice of recent temporal reasoning benchmarks like \cite{zhou2019going, wang-zhao-2024-tram}, using a multiple-choice format for reliable assessment. A more ideal setting would be evaluating in generative dialogue settings. However, generative QA evaluations pose significant challenges due to variations in temporal expressions (e.g., "January 1st" vs. "01/01"), making exact match and F1 token score unreliable. Thus, we prioritize benchmark reliability and accuracy over a more ambitious generative QA setting. We leave the extension of the current benchmark to generative dialogue settings as future work

\noindent\textbf{Open-sourced Models.} Our current experiments mainly focus on close-sourced models. However, as we have pointed out in our Experiment section $\S$ \ref{sec:models}, we found most open-source LLMs cannot handle the long dialogue inputs from LoCoMo, for example only about 10\% of dialogues can be fed into Llama-3-70B. And even for those shorter dialogues that can be fed into Llama-3-70B, we notice that the model gets lost and fails to follow instructions, even failing to generate in the desired format. We therefore consider the adaptation of our framework to open-sourced models as future work. Note that we will release our code and dataset for reproducibility.

% Additionally, we observe that the LoCoMo dialogues contain some temporal ambiguities that can affect reasoning over related events. For instance, if a speaker mentions on a Tuesday that an event occurred "last Monday," it is unclear whether the event happened the previous day or on the Monday of the prior week. However, in real applications, LLM agents could potentially resolve such ambiguities by asking follow-up clarifying questions.